\section{GIỚI THIỆU BÀI TOÁN}
\subsection{Đặt vấn đề}
Trong các hệ thống hỗ trợ lái xe và giao thông thông minh, biển báo cung cấp những thông tin trực tiếp về giới hạn tốc độ, hướng di chuyển và các tình huống cần cảnh báo. Khả năng nhận dạng đúng biển báo phụ thuộc trước hết vào chất lượng hình ảnh mà hệ thống tiếp nhận. Khi điều kiện chiếu sáng suy giảm, những đặc trưng quan trọng như đường viền, chữ số và mũi tên có thể trở nên khó phân biệt, ngay cả khi vị trí của biển báo đã được xác định.

Một hướng giải quyết thường được cân nhắc là tăng cường ảnh trước khi đưa vào mô hình nhận dạng. Tuy nhiên, tăng độ sáng không đồng nghĩa với phục hồi đầy đủ thông tin. Một phép biến đổi có thể làm vùng tối dễ quan sát hơn nhưng đồng thời khuếch đại nhiễu, thay đổi màu sắc hoặc làm nhòe những chi tiết nhỏ. Đối với biển báo giao thông, chỉ một nét bị mất trong chữ số hoặc một hướng mũi tên bị biến dạng cũng có thể dẫn đến dự đoán sai lớp.

Từ đó, bài toán được đặt ra là: \textit{các phương pháp xử lý ảnh cải thiện chất lượng ảnh thiếu sáng đến mức nào, và sự cải thiện đó có thực sự giúp mô hình nhận dạng chính xác hơn hay không?} Đây là câu hỏi cần được trả lời bằng cả thước đo chất lượng ảnh lẫn kết quả nhận dạng trên cùng một điều kiện thực nghiệm.

\textit{Thứ nhất, suy giảm độ sáng đi kèm với suy giảm tỷ lệ tín hiệu trên nhiễu.} Vùng tối có biên độ tín hiệu nhỏ, do đó nhiễu có thể chiếm tỷ lệ đáng kể so với thông tin hữu ích. Nếu chỉ áp dụng phép tăng sáng, phần nhiễu cũng được nâng lên cùng với cấu trúc của biển báo. Vì vậy, đánh giá một phương pháp trên ảnh chỉ làm tối bằng Gamma sẽ chưa phản ánh đầy đủ trường hợp ảnh đồng thời chứa nhiễu.

\textit{Thứ hai, chất lượng ảnh và chất lượng đặc trưng nhận dạng không hoàn toàn trùng nhau.} PSNR phản ánh sai số cường độ điểm ảnh, còn SSIM xem xét độ tương đồng về cấu trúc cục bộ. Trong khi đó, mô hình nhận dạng học những đặc trưng phục vụ phân lớp. Một ảnh có sai số trung bình thấp vẫn có thể làm mất nét của chữ số; ngược lại, một ảnh khác ảnh tham chiếu về độ sáng vẫn có thể giữ được hình dạng cần thiết để nhận dạng.

\textit{Thứ ba, so sánh thuật toán cần kiểm soát các yếu tố ngoài xử lý ảnh.} Nếu thay đổi kiến trúc mô hình hoặc huấn luyện lại riêng cho từng phương pháp tăng cường, rất khó tách biệt phần cải thiện đến từ ảnh và phần cải thiện đến từ mô hình. Bài tập lớn vì thế sử dụng duy nhất kiến trúc ResNet-18, huấn luyện trên ảnh gốc và giữ cố định trọng số khi so sánh các phép xử lý.

\subsection{Yêu cầu và phạm vi nghiên cứu}
Đề bài yêu cầu khảo sát ảnh hưởng của xử lý ảnh đến hiệu năng nhận dạng biển báo giao thông trong điều kiện thiếu sáng trên bộ dữ liệu GTSRB. Ảnh suy giảm được mô phỏng bằng Gamma kết hợp nhiễu Gaussian và Poisson ở nhiều mức độ. Các phương pháp tăng cường cần so sánh gồm hiệu chỉnh Gamma, cân bằng histogram (HE), cân bằng histogram cục bộ có giới hạn tương phản (CLAHE), Retinex, Zero-DCE, Retinexformer và phương pháp dựa trên diffusion.

Nghiên cứu còn so sánh Gaussian, bilateral và Non-local Means khi dùng đơn lẻ, cùng hai thứ tự phối hợp với Gamma và CLAHE. Độ mạnh bộ lọc được chọn trên validation và giữ giống nhau trong từng cặp thứ tự. Chất lượng ảnh được đánh giá bằng PSNR và SSIM. Hiệu năng nhận dạng được đo bằng Accuracy và F1-score, trong đó Macro-F1 là tiêu chí chính để tránh đánh giá thiên lệch về các lớp có nhiều ảnh. Kết quả được dùng để phân tích tương quan giữa chất lượng ảnh và khả năng nhận dạng, từ đó lựa chọn quy trình tiền xử lý phù hợp.

Trong phạm vi triển khai, đầu vào là vùng ảnh đã chứa biển báo, không phải toàn cảnh giao thông. Hệ thống giải quyết bài toán phân loại vào một trong 43 lớp của GTSRB; chưa bao gồm phát hiện vị trí biển báo, theo dõi đối tượng hoặc xử lý video. Điều kiện thiếu sáng được tạo có kiểm soát từ ảnh gốc, do đó có ảnh tham chiếu tương ứng để tính PSNR và SSIM.

\subsection{Đầu vào và đầu ra của hệ thống}
\textbf{Đầu vào:} ảnh màu RGB sau khi cắt vùng quan tâm và đưa về kích thước $64\times64$ pixel. Giá trị cường độ được chuẩn hóa về đoạn $[0,1]$. Từ ảnh gốc $X$, mô hình suy giảm tạo ảnh thiếu sáng $Y$ với một mức tối và một hạt giống sinh nhiễu xác định.

\textbf{Đầu ra của khối xử lý ảnh:} ảnh tăng cường $\widehat{X}=E(Y)$ có cùng kích thước, thứ tự kênh màu và miền giá trị với đầu vào. Hàm $E$ chỉ nhận ảnh thiếu sáng; ảnh tham chiếu và nhãn lớp không được đưa vào thuật toán tăng cường.

\textbf{Đầu ra của khối nhận dạng:} vector 43 giá trị dự đoán và nhãn lớp có giá trị lớn nhất. Trước khi đi qua ResNet-18, ảnh được nội suy lên $224\times224$ và chuẩn hóa theo ImageNet. Việc đổi kích thước này chỉ phục vụ mô hình nhận dạng; các phép tăng cường và đo chất lượng ảnh vẫn thực hiện ở $64\times64$.

\textbf{Đầu ra nghiên cứu:} bảng so sánh PSNR, SSIM, Accuracy, Macro-F1, thời gian xử lý; hình ảnh trực quan; phân tích ablation và tương quan; cùng một phương án tiền xử lý được chọn trên validation trước khi đánh giá test.

\subsection{Các hướng tiếp cận và giải pháp khảo sát}
Nhóm phương pháp truyền thống tác động trực tiếp lên cường độ hoặc phân bố mức xám. Gamma điều chỉnh phi tuyến từng điểm ảnh; HE phân phối lại histogram toàn cục; CLAHE xử lý theo vùng; Retinex ước lượng quan hệ giữa thành phần chiếu sáng và phản xạ. Những phương pháp này có ưu điểm dễ giải thích, ít tham số và chi phí tính toán tương đối thấp.

Nhóm phương pháp học sâu sử dụng thông tin đã học để điều chỉnh hoặc phục hồi ảnh. Zero-DCE ước lượng đường cong tăng sáng, Retinexformer kết hợp biểu diễn chiếu sáng với Transformer, còn phương pháp diffusion được khảo sát thực hiện phục hồi có điều kiện trong miền wavelet. Dự án sử dụng trọng số công bố của tác giả, không huấn luyện lại các bộ tăng cường trên GTSRB~\cite{zerodce}, \cite{retinexformer}, \cite{diffusion}.

Giải pháp của bài tập lớn không giả định một nhóm phương pháp luôn tốt hơn nhóm còn lại. Tất cả được đánh giá trên cùng ảnh suy giảm, cùng tập mẫu và cùng các checkpoint ResNet-18. Mốc đối chứng là \term{Identity}, tức ảnh tối giữ nguyên. Ảnh gốc \term{Clean} được dùng làm tham chiếu và kiểm tra hiệu năng nhận dạng khi chưa áp dụng mô hình suy giảm.

\subsection{Đóng góp triển khai và xây dựng quy trình thực nghiệm}
\label{sec:contributions}
Nhóm xây dựng quy trình chuẩn bị dữ liệu từ các archive GTSRB chính thức, gồm đọc ảnh PPM và annotations, chuyển ảnh sang RGB, cắt vùng biển báo theo ROI với padding 5\% và nội suy song tuyến tính về kích thước $64\times64$. Các bước này vận dụng kiến thức về biểu diễn ảnh, hệ tọa độ, phép cắt và lấy mẫu lại; đồng thời tạo đầu vào thống nhất cho tất cả phương pháp tăng cường và bộ phân loại. Chi tiết cài đặt và ý nghĩa của các lựa chọn được trình bày trong phần xây dựng tập dữ liệu.

Trong quá trình mô phỏng suy giảm và tăng cường, ảnh được biểu diễn bằng số thực trong miền $[0,1]$; nhóm tăng sáng lưu kết quả trung gian bằng NumPy, còn khảo sát khử nhiễu/kết hợp truyền trực tiếp mảng float32 trong bộ nhớ để hạn chế lượng tử hóa và I/O bổ sung. Riêng HE, CLAHE và NLM có bước chuyển sang ảnh 8 bit theo cài đặt OpenCV, được mô tả rõ trong phần phương pháp. Việc thống nhất miền giá trị, thứ tự kênh và kích thước ảnh giúp tránh sai khác ngoài thuật toán khi so sánh PSNR, SSIM và hiệu năng phân loại.

Nhóm kết hợp quy trình xử lý ảnh với kiểm tra trùng lặp, chia dữ liệu theo track và lưu checksum cùng thông tin nguồn để hỗ trợ tái lập thí nghiệm. Giá trị của phần triển khai nằm ở khả năng kiểm soát, giải thích và kiểm chứng toàn bộ luồng dữ liệu; thao tác tải, giải nén và gom annotations là các bước hỗ trợ. Đây là đóng góp về triển khai và thiết kế thực nghiệm, không phải đề xuất thuật toán xử lý ảnh mới. Padding, kích thước crop và phép nội suy được trình bày là các quyết định thiết kế; chưa có thí nghiệm so sánh riêng để khẳng định chúng là cấu hình tối ưu.

\subsection{Ý nghĩa và bố cục báo cáo}
Về kiến thức xử lý ảnh, đề tài kết nối phép biến đổi điểm, histogram, lọc Gaussian, mô hình nhiễu và các chỉ số chất lượng trong một quy trình thống nhất. Về thực nghiệm, đề tài làm rõ sự đánh đổi giữa độ sáng, mức nhiễu, tính toàn vẹn của chi tiết và chi phí suy luận. ResNet-18 đóng vai trò công cụ kiểm chứng thông tin được giữ lại sau xử lý, thay vì là đối tượng tối ưu chính của báo cáo.

Các phần tiếp theo lần lượt trình bày dữ liệu và mô hình suy giảm; cơ sở và cài đặt các phương pháp tăng cường; thiết lập cùng kết quả thực nghiệm; phân tích tương quan và ảnh hưởng của nhiễu; quy trình tiền xử lý được đề xuất; cuối cùng là kết luận và giới hạn của nghiên cứu.
